\section{Poisson Processes}

  \subsection{Exponential Distribution}

    Let us do some review. The \textbf{exponential distirbution} of rate $\lambda$ is a \hyperref[prob-def:random-variable]{random variable} $T \sim \mathrm{Exponential}(\lambda)$ with CDF 

      \begin{equation}
        F_T (t) = \mathbb{P}(T \leq t) = 1 - e^{-\lambda t}
      \end{equation}

    and the PDF 

      \begin{equation}
        f_T (t) = \begin{cases} \lambda e^{-\lambda t} & t \geq 0 \\ 0 & t < 0 \end{cases}
      \end{equation}

    We have 

      \begin{equation}
        \mathbb{E}[T] = \frac{1}{\lambda}, \;\; \mathrm{Var}(T) = \frac{1}{\lambda^2}
      \end{equation}

    \begin{lemma}[Memoryless Property]
      \label{thm:memoryless-property}
      The Exp$(\lambda)$ distribution has the property that for all $t, s \geq 0$, 

        \begin{equation}
          \mathbb{P}(W > t + s \; | \; W > t) = \mathbb{P}(W > s)
        \end{equation}

      which is called the \textit{memoryless property}. We can interpret this in the following way. Let $W$ be the time you have to wait for the first arrival. Given that you already waited $t$ units of time, the probability that you have the wait $s$ additional units of time is just the probability that you wait at least $s$ from the beginning. That is, knowing that $t$ units of time have elapsed does not affect the distribution of the remaining waiting time. 
    \end{lemma}

    \begin{theorem}[Characterization of the Exponential Distribution]
      \label{thm:characterization-of-the-exponential-distribution}
      Let $W$ be a continuously distributed \hyperref[prob-def:random-variable]{random variable}. Then $W \sim$ Exp$(\lambda)$ for some $\lambda > 0$ if and only if $W$ satisfies the \hyperref[thm:memoryless-property]{memoryless property}. 
    \end{theorem}

    \begin{theorem}[Minimum of Independent Exponentials]
      \label{thm:minimum-of-independent-exponentials}
      Let $T_i \sim \mathrm{Exponential}(\lambda_i)$ for $i = 1, \ldots n$. Then, 

        \begin{equation}
          \min\{T_1, \ldots, T_n\} \sim \mathrm{Exponential}(\lambda_1 + \ldots + \lambda_n)
        \end{equation}

      and the \hyperref[prob-def:random-variable]{random variable} $I$ which takes the index of $\min\{T_1, \ldots, T_n\}$ has the PMF 

        \begin{equation}
          \mathbb{P}(I = i) = \frac{\lambda_i}{\lambda_1 + \ldots + \lambda_n}
        \end{equation}
    \end{theorem}

  \subsection{Defining the Poisson Process}

    We first describe a limiting behavior of \hyperref[prob-def:binomial-random-variable]{binomial random variables}. 

    \begin{theorem}[Poisson Limit Theorem]
      \label{thm:poisson-limit-theorem}
      Let $X_n \sim \mathrm{Bernoulli}(n, p_n)$, where $\{p_n\}_{n \in \mathbb{N}}$ is a sequence of reals in $[0, 1]$ such that 

        \begin{equation}
          \lim_{n \rightarrow \infty} n p_n = \lambda
        \end{equation}

      Letting $Y \sim \mathrm{Poisson}(\lambda)$

        \begin{equation}
          X_n \xrightarrow{D} Y
        \end{equation}

      That is, the CDFs, and since this is a discrete distribution, the PMFs, converge. 
    \end{theorem}
    \begin{proof}
      We will show that $\lim_{n \rightarrow \infty} \mathbb{P}(X_n = k) = \mathbb{P}(Y = k)$, which shows that the CDFs converge and therefore \hyperref[prob-def:convergence-in-distribution]{convergence in distribution}. 
      \begin{align*}
        \lim_{n \rightarrow \infty} \mathbb{P}(X_n = k) & = \lim_{n \rightarrow \infty} \binom{n}{k} p_n^k (1 - p_n)^k \\
        & = \lim_{n \rightarrow \infty} \frac{n (n - 1) \ldots (n - k + 1)}{k!} \bigg(\frac{\lambda}{n}\bigg)^k \bigg( 1 - \frac{\lambda}{n} \bigg)^{n - k} \\
        & = \lim_{n \rightarrow \infty} \frac{n^k + O(n^{k-1})}{k!} \frac{\lambda^k}{n^k} \bigg( 1 - \frac{\lambda}{n} \bigg)^{n - k} \\
        & = \lim_{n \rightarrow \infty} \frac{\lambda^k}{k!} \bigg( 1 - \frac{\lambda}{n} \bigg)^{n - k} \\
        & = \frac{\lambda^k}{k!} \lim_{n \rightarrow \infty} \bigg( 1 - \frac{\lambda}{n} \bigg)^n \; \lim_{n \rightarrow \infty} \bigg( 1 - \frac{\lambda}{n} \bigg)^{-k} \\
        & = \frac{\lambda^k}{k!} \, e^{\lambda} \, 1 = \frac{\lambda^k e^\lambda}{k!}
      \end{align*}
    \end{proof}

    Note that this is different from CLT because in CLT, we just assume that the $p_n$'s are constant and take the limiting behavior of $X_n \sim \mathrm{Bernoulli}(n, p)$ as $n \rightarrow \infty$. 

    This result justifies the following model. A Poisson Arrival Process with rate $\lambda > 0$ on the interval $[0, \infty)$ is a model for the occurrence of some events which may have at any time. We can interpret the process as a collection of random points in $[0, \infty)$ which are the times at which the arrivals occur. Suppose that we would like to model the arrival of events that happen completely at random at a rate $\lambda$ per unit time. At time $t = 0$, we have no arrivals yet, so $N(0) = 0$. Let us fix some $T$, and now divide $[0, T)$ into $n$ tiny subintervals of length $\delta$. 

    \begin{center}
      %\includegraphics[scale=0.3]{}
    \end{center}

    Assume that in each time slot, we assign a $X_k \sim \mathrm{Bernoulli}(\lambda \delta)$ \hyperref[prob-def:random-variable]{random variable} that determines whether there was an arrival within the interval $((k-1)\delta, k\delta]$. So with probability $\lambda\delta$, there will be an arrival within it, and as the time interval gets smaller, this probability also gets smaller too. Since every $n$ subinterval is $\mathrm{Bernoulli}(\lambda \delta)$, the number of arrivals in the interval $[0, T)$, defined as the random variable $N_n (T)$, is 

      \begin{equation}
        N_n (T) \sim \mathrm{Binomial}(n, \lambda \delta) = \mathrm{Binomial}\big(n, \frac{\lambda T}{n}\big)
      \end{equation}

    As we increase the $n$ (equivalently, decrease $\delta$), we divide $[0, T)$ into smaller and smaller subintervals, resulting in finer and finer $N_n(T)$ Binomial distributions. Since $n p_n = n \frac{\lambda T}{n} = \lambda T$ is finite, we can invoke the \hyperref[thm:poisson-limit-theorem]{Poisson limit theorem} and say 

      \begin{equation}
        N_n (T) \xrightarrow{D} \mathrm{Poisson}(\lambda T)
      \end{equation}

    Note that the starting point $0$ does not matter, and this works for any interval of length $T$. Therefore, we can model the arrival times on any interval of length $T$ as a $\mathrm{Poisson}(\lambda T)$ random variable. 

    \begin{definition}[Poisson Process]
      \label{def:poisson-process}
      Let $\lambda > 0$ be fixed, representing the rate of arrival in some unit time. The stochastic counting process $\{N(t)\}_{t \geq 0}$, where $N(t)$ represents the number of arrivals by time $t$, is called a \textbf{Poisson process} with rate $\lambda$ if 
      \begin{enumerate}
        \item $N(0) = 0$ 

        \item The number of arrivals in any interval of length $s > 0$ is $N(t + s) - N(t) \sim \mathrm{Poisson}(\lambda s)$ 

        \item $N(Tt)$ has independent increments, i.e. if $t_0 < t_1 < \ldots, < t_n$, then 

          \begin{equation}
            N(t_1) - N(t_0), \ldots, N(t_n) - N(t_{n-1})
          \end{equation}

        are independent. 
      \end{enumerate}
    \end{definition}

  \subsection{Constructing the Poisson Process}

    Now we have modeled this process using \hyperref[prob-def:random-variable]{random variables} $N(t)$ that counts the number of arrivals up to time $t$. Now, we can interpret it using random variables that represent the \textit{time} in which they arrive. 

    \begin{definition}[Arrival and Inter-Arrival Times]
      \label{def:arrival-and-inter-arrival-times}
      Set $T_0 = 0$. The arrival times are \hyperref[prob-def:random-variable]{random variables} $0 < T_1 < T_2 < T_3 < \ldots$ such that the inter-arrival waiting times

        \begin{equation}
          \tau_k = T_k - T_{k-1}, \;\;\; k \geq 0
        \end{equation}

      have the property that $\{\tau_k\}_{k=1}^\infty$ are independent $\mathrm{Exp}(\lambda)$ random variables. Define 

        \begin{equation}
          N(s) \coloneqq \max\{k \mid T_k \leq s \}
        \end{equation}
    \end{definition}

    Now we prove that this process is equivalent to the \hyperref[def:poisson-process]{Poisson process} defined before. 

    \begin{theorem}[Equivalent Interpretations]
      \label{thm:equivalent-interpretations}
      Let $\{T_n\}$ be defined as above and $N(s) \coloneqq \max\{k \mid T_k \leq s \}$. Then, 

      \begin{enumerate}
        \item $N(0) = 0$
        \item $N(s) \sim \mathrm{Poisson}(\lambda s)$ 
        \item $N(t + s) - N(t) \sim \mathrm{Poisson}(\lambda s)$ independent of $N(r)$ for $0 \leq r \leq s$. 
        \item $N(t)$ has independent incremements. 
      \end{enumerate}

      $N(s) \coloneqq \max\{k \mid T_k \leq s \}$ is a Poisson distribution with mean $\lambda s$. 
    \end{theorem}

    \begin{proof}
      From the second perspective, the statement $N(I) \sim \mathrm{Poisson}(\lambda |I|)$ for an interval $I$ means that

        \begin{equation}
          \mathbb{P}(N(I) = m) = e^{-\lambda |I|} \frac{(\lambda |I|)^m}{m!}, \qquad m = 0, 1, 2, \ldots
        \end{equation}

      where $|I|$ is the length of the interval $I$. From the first perspective, notice that

        \begin{equation}
          T_k = \tau_1 + \tau_2 + \ldots + \tau_k
        \end{equation}

      so that the $k$th arrival time $T_k$ is a sum of $k$ independent $\mathrm{Exp}(\lambda)$ \hyperref[prob-def:random-variable]{random variables}. Thus

        \begin{equation}
          T_k \sim \mathrm{Gamma}(k, \lambda)
        \end{equation}

      and therefore has density

        \begin{equation}
          \lambda e^{-\lambda t} \frac{(\lambda t)^{k-1}}{(k-1)!}, \qquad t > 0
        \end{equation}

      Note that the arrival times $T_i$ are not independent of each other, but the inter-arrival times $\tau_i$ are indeed independent.
    \end{proof}

  \subsection{The Poisson Process on a Finite Horizon}

    We can slightly modify the construction above to create a \hyperref[def:poisson-process]{Poisson process} over some finite time horizon $[0, L]$. Again, there are two equivalent ways to do this.

    \begin{enumerate}
      \item Starting with independent $\mathrm{Exp}(\lambda)$ \hyperref[prob-def:random-variable]{random variables} $\tau_1, \tau_2, \ldots$, we define

        \begin{equation}
          T_k = \sum_{i=1}^k \tau_i
        \end{equation}

      and stop once $T_k > L$.

      \item We let $N \sim \mathrm{Poisson}(\lambda L)$, since we are only working on the finite interval $[0, L]$. Given $N = n$, let $U_1, U_2, \ldots, U_n \sim \mathrm{Uniform}([0, L])$ be independent. These define the arrival times, and ordering them gives

        \begin{equation}
          T_k = U_{(k)}, \qquad k = 1, 2, \ldots, N
        \end{equation}

      where $U_{(k)}$ is the $k$th order statistic, so that $T_1 = \min(U_1, \ldots, U_N)$.
    \end{enumerate}
